\begin{quote}
Usar como guion para construir el deck (Beamer / PowerPoint / Google Slides). $\sim$16--22 slides.
\end{quote}
\paragraph{Bloque A --- Apertura}\begin{enumerate}
\item Portada (CADi, sesión 1, Juliho Castillo, EIC)
\item Mapa 40\,h (24 aula + 16 independientes) y producto del día (E1)
\item Resultados de aprendizaje de la sesión
\item Normas de taller e integridad (\ceddie)
\end{enumerate}
\paragraph{Bloque B --- Conceptos}\begin{enumerate}
\item De chat a agente
\item Anatomía de un agente
\item Preview de orquestación (detalle en sesión 2)
\item Human-in-the-loop (HITL)
\item Roles STEM (planificador, redactor, revisor, verificador)
\end{enumerate}
\paragraph{Bloque C --- Integridad}\begin{enumerate}
\item Riesgos (alucinación, citas, sesgo, datos)
\item Ecuaciones, unidades y citas
\item Transparencia y agent log
\end{enumerate}
\paragraph{Bloque D --- Demo y taller}\begin{enumerate}
\item Caso: idea vaga vs.\ flujo con HITL
\item Anti-patrones
\item Instrucciones de la actividad (E1)
\item Rúbrica rápida
\end{enumerate}
\paragraph{Bloque E --- Cierre}\begin{enumerate}
\item Takeaways
\item TI bloque A ($\sim$1\,h)
\item Prep sesión 2
\item Cierre / gracias
\end{enumerate}

\subsubsection*{Notas de diseño}
\begin{itemize}
\item Máximo 1 idea fuerte por slide.
\item Evitar paredes de texto; mover detalle a la guía del facilitador.
\item Mantener marcadores CEDDIE / verificar política institucional donde corresponda.
\item Incluir recordatorio de las \textbf{16\,h} de trabajo independiente cuando haya encargo TI.
\end{itemize}
